% ------------------------------------------------------------
% Preamble for 「銀河力学で学ぶ数値計算法」
% Style guide v1.1: A4, derivation-centered, restrained emphasis.
% Compile with XeLaTeX.
% ------------------------------------------------------------
\usepackage[a4paper,inner=27mm,outer=22mm,top=22mm,bottom=25mm,
  headheight=14pt,headsep=6mm]{geometry}

\usepackage{amsmath,amssymb,mathtools,bm}
\usepackage{graphicx,float,caption}
\captionsetup{font=small,labelfont=bf,labelsep=quad}
\graphicspath{{figures/}{./}}

\usepackage{fontspec}
\usepackage{xeCJK}
\IfFontExistsTF{Latin Modern Roman}{\setmainfont{Latin Modern Roman}}{\setmainfont{Liberation Serif}}
\IfFontExistsTF{Latin Modern Sans}{\setsansfont{Latin Modern Sans}}{\setsansfont{Liberation Sans}}
\IfFontExistsTF{DejaVu Sans Mono}{\setmonofont{DejaVu Sans Mono}[Scale=0.91]}{}
\setCJKmainfont{Noto Serif CJK JP}[ItalicFont={Noto Serif CJK JP}]
\setCJKsansfont{Noto Sans CJK JP}
\setCJKmonofont{Noto Sans Mono CJK JP}[Scale=0.90]
\xeCJKsetup{CJKmath=true}
\XeTeXlinebreaklocale "ja"
\XeTeXlinebreakskip = 0pt plus 0.1pt minus 0.01pt

\usepackage{xcolor,enumitem,array,booktabs,tabularx,xurl}
\usepackage{tikz}
\usetikzlibrary{arrows.meta,calc,positioning,patterns,decorations.pathreplacing,matrix}
\usepackage[most]{tcolorbox}
\usepackage{fancyvrb}
% Code blocks in Listing and Notebook excerpts use line numbers for discussion in class.
\fvset{numbers=left,numbersep=5pt,firstnumber=1,stepnumber=1}
\renewcommand{\theFancyVerbLine}{\sffamily\scriptsize\color{graytext}\arabic{FancyVerbLine}}
\usepackage{titlesec,fancyhdr}

\setlength{\parindent}{1em}
\setlength{\parskip}{0.15em}
\linespread{1.08}
\raggedbottom
\emergencystretch=1.5em
\widowpenalty=10000
\clubpenalty=10000

% --- Color system inherited from the simulation text, with restrained accents.
\definecolor{navy}{HTML}{0D1B3E}
\definecolor{cyan}{HTML}{1F6390}
\definecolor{cyanbg}{HTML}{E7F0F7}
\definecolor{turquoise}{HTML}{3A8F9B}
\definecolor{turquoisebg}{HTML}{EDF7F8}
\definecolor{amber}{HTML}{A8680B}
\definecolor{amberbg}{HTML}{FFF2CF}
\definecolor{rule}{HTML}{C9D2E3}
\definecolor{graytext}{HTML}{555555}
\definecolor{listingbg}{HTML}{F7F8FA}
\definecolor{listingframe}{HTML}{B8C0CC}

% --- Chapter-opening learning box: only full-width pedagogical box used routinely.
\newtcolorbox{learningbox}[1][本章で学ぶこと]{
  enhanced,breakable,
  colback=cyanbg,colframe=cyan,boxrule=0.65pt,arc=1.5pt,
  left=3mm,right=3mm,top=2.2mm,bottom=2.2mm,
  colbacktitle=cyan,coltitle=white,fonttitle=\sffamily\bfseries,
  title={#1},before skip=5pt,after skip=8pt}

% --- Kouya-inspired thin box: only for definitions/theorems with independent meaning.
\newtcolorbox{definitionbox}[2][]{
  enhanced,breakable,
  colback=white,colframe=turquoise,boxrule=0.6pt,arc=0pt,
  left=3mm,right=3mm,top=5.5mm,bottom=2.2mm,
  before skip=7pt,after skip=7pt,
  overlay unbroken and first={
    \node[anchor=north west,fill=turquoise,text=white,
      font=\sffamily\bfseries\small,inner xsep=2.2mm,inner ysep=0.8mm]
      at (frame.north west) {#2};
  },#1}

% --- Scoped supplements: white background + left rule only.
\newtcolorbox{notebookblock}[1]{
  enhanced,breakable,frame hidden,boxrule=0pt,arc=0pt,
  colback=white,borderline west={2.2pt}{0pt}{turquoise},
  left=3.2mm,right=3mm,top=5.3mm,bottom=2mm,
  before skip=8pt,after skip=8pt,
  title={Notebook #1},
  coltitle=turquoise,colbacktitle=white,
  fonttitle=\sffamily\bfseries\small,
  boxed title style={boxrule=0pt,frame hidden,colback=white},
  attach boxed title to top left={xshift=0mm,yshift=-0.5mm}}

\newtcolorbox{cautionblock}[1]{
  enhanced,frame hidden,boxrule=0pt,arc=0pt,
  colback=white,borderline west={2.2pt}{0pt}{amber},
  left=3.2mm,right=3mm,top=5.3mm,bottom=2mm,
  before skip=8pt,after skip=8pt,
  title={数値計算上の注意｜#1},
  coltitle=amber,colbacktitle=white,
  fonttitle=\sffamily\bfseries\small,
  boxed title style={boxrule=0pt,frame hidden,colback=white},
  attach boxed title to top left={xshift=0mm,yshift=-0.5mm}}

% 本編の理解には必須でない関連手法を、補足として紹介する。
% 白背景と青い左線のみ。長い補足は必要に応じてページをまたいで表示する。
\newtcolorbox{topicblock}[1]{
  enhanced,breakable,
  frame hidden,boxrule=0pt,arc=0pt,
  colback=white,borderline west={2.2pt}{0pt}{cyan},
  left=3.2mm,right=3mm,top=5.3mm,bottom=2mm,
  before skip=8pt,after skip=8pt,
  before upper={\setlength{\parindent}{1em}\indent},
  title={補足｜#1},
  coltitle=cyan,colbacktitle=white,
  fonttitle=\sffamily\bfseries\small,
  boxed title style={boxrule=0pt,frame hidden,colback=white},
  attach boxed title to top left={xshift=0mm,yshift=-0.5mm}}

% Short Python code bridging a derived formula and the complete Notebook.
\newtcolorbox{listingbox}[2][]{
  enhanced,
  colback=listingbg,colframe=listingframe,boxrule=0.45pt,arc=0pt,
  left=2.8mm,right=2.8mm,top=6mm,bottom=2.5mm,
  before skip=8pt,after skip=8pt,
  overlay unbroken and first={
    \node[anchor=north west,fill=white,text=navy,
      font=\sffamily\bfseries\small,inner xsep=2.0mm,inner ysep=0.7mm]
      at ([xshift=1.2mm]frame.north west) {Listing #2};
  },#1}


% Compact source syntax; rendered as the scoped Notebook block.
\newcommand{\notebooktag}[2]{%
  \begin{notebookblock}{#1}#2\end{notebookblock}}

\titleformat{\chapter}{\LARGE\bfseries\sffamily\color{navy}}{第\thechapter 章}{0.8em}{}
\titleformat{\section}{\large\bfseries\sffamily\color{navy}}{\thesection}{0.5em}{}
\titleformat{\subsection}{\normalsize\bfseries\sffamily\color{navy}}{}{0pt}{}
\titlespacing*{\chapter}{0pt}{0pt}{12pt}
\titlespacing*{\section}{0pt}{1.1em}{0.4em}
\titlespacing*{\subsection}{0pt}{0.75em}{0.25em}

\setlist[itemize]{leftmargin=1.7em,itemsep=0.15em,topsep=0.25em}
\setlist[enumerate]{leftmargin=2em,itemsep=0.15em,topsep=0.25em}
\setcounter{tocdepth}{1}
\setcounter{secnumdepth}{1}

\pagestyle{fancy}
\fancyhf{}
\fancyhead[L]{\scriptsize\sffamily 銀河力学で学ぶ数値計算法}
\fancyhead[R]{\scriptsize\sffamily\nouppercase{\leftmark}}
\makeatletter
\renewcommand{\chaptermark}[1]{\markboth{第\thechapter 章\enspace #1}{}}
\makeatother
\fancyfoot[R]{\small\thepage}
\fancypagestyle{plain}{
  \fancyhf{}\fancyfoot[R]{\small\thepage}\renewcommand{\headrulewidth}{0pt}}

\newcommand{\unit}[1]{\,\mathrm{#1}}
\newcommand{\dd}{\mathrm{d}}
\renewcommand{\contentsname}{目次}
\renewcommand{\figurename}{図}
\renewcommand{\tablename}{表}
\renewcommand{\bibname}{参考図書}
\renewcommand{\thefigure}{\thechapter.\arabic{figure}}
\renewcommand{\thetable}{\thechapter.\arabic{table}}
\renewcommand{\theequation}{\thechapter.\arabic{equation}}

\renewcommand{\topfraction}{0.95}
\renewcommand{\bottomfraction}{0.85}
\renewcommand{\textfraction}{0.05}
\renewcommand{\floatpagefraction}{0.85}
\setcounter{topnumber}{3}
\setcounter{totalnumber}{5}
\setlength{\textfloatsep}{10pt plus 3pt minus 2pt}
\setlength{\intextsep}{10pt plus 3pt minus 2pt}

\usepackage[unicode,colorlinks=true,linkcolor=cyan,citecolor=cyan,urlcolor=cyan]{hyperref}
